% paper/sections/04_time_integration.tex
%
% §4  CLS 移流の離散化概要
%     WENO5 の完全な定式化は付録~\ref{app:weno5} に移動した．
%     本章では移流スキーム選択の動機と TVD-RK3 の接続のみを扱う．
%
\section{CLS 移流の離散化概要}
\label{sec:time}

本章では，第~\ref{sec:levelset}章で導入した CLS 移流方程式
\begin{equation}
  \frac{\partial\psi}{\partial t} + \nabla\cdot(\psi\bm{u}) = 0
  \label{eq:cls_advection_ref}
\end{equation}
（式~\eqref{eq:cls_advection}の再掲）を数値的に解くための\textbf{空間離散化スキームの選択根拠}を示す．
CCD による微分計算（第~\ref{sec:CCD}章）および
散逸的 CCD フィルタ（§\ref{sec:dissipative_ccd}）の設計を踏まえ，
本章はそれらと CLS 移流がどのように役割分担するかを説明する．

\subsection{スキーム選択の動機：なぜ WENO5 が必要か}
\label{sec:advection_motivation}

\subsubsection{気液界面移流の困難}

$\psi$ は界面付近で $0$ から $1$ へと急激に変化するが，この勾配の強さは
物理的実体（物性の不連続）を反映したものであって，数値拡散でなだらかにして良いものではない．

\textbf{線形スキームの限界：}
通常の線形スキーム（例：CCD 直接適用）は高精度だが，
$\psi$ のような急峻なプロファイルを移流すると
\textbf{Gibbs 振動}（$\psi < 0$ や $\psi > 1$ の非物理的値）が発生する．
これが界面の「にじみ」（第~\ref{sec:intro}章~§\ref{sec:failure_modes}）の主因の一つである．

\textbf{CCD との役割分担：}
CCD は\textbf{滑らかな場}の高精度微分評価（曲率 $\kappa$，速度勾配 $\nabla\bm{u}$，
圧力勾配 $\nabla p$）に最適化されており，急峻な勾配を持つ $\psi$ の\textbf{移流}には
設計上適合しない（中心差分ベースの振動問題）．
一方，散逸的 CCD フィルタ（§\ref{sec:dissipative_ccd}）は CCD 微分場の
高周波ノイズを除去するが，$\psi$ の非線形移流そのものを安定化するには至らない．

したがって，CLS 移流には\textbf{不連続・急峻勾配に対応した非線形スキーム}が別途必要である．

\subsubsection{WENO5 の採用}

本稿では，CLS 移流に 5 次精度の Weighted ENO（WENO5）スキーム \cite{JiangShu1996} を採用する．
WENO5 は\textbf{滑らかさ指標（smoothness indicator）$\beta_k$} に基づく非線形重みを用い，
\begin{itemize}
  \item \textbf{滑らかな領域：} 3 つの候補ステンシルを最適に凸結合し $O(h^5)$ 精度を達成
  \item \textbf{不連続・急峻勾配付近：} 不連続を含む候補ステンシルの重みが自動的にゼロとなり
        振動を抑制（ENO 的挙動）
\end{itemize}
する自己適応型スキームである．

\textbf{保存形との整合性：}
式~\eqref{eq:cls_advection_ref} は保存形（発散形）で記述されており，
WENO5 もフラックス形式（$\hat{F}_{i+1/2}$）で実装されるため，
離散的な質量保存が自動的に成立する（CLS 法の主要な利点，§\ref{sec:why_cls}）．

\medskip
WENO5 の完全な定式化（Lax--Friedrichs フラックス分割，
滑らかさ指標 $\beta_k$，再構成係数，境界ゴーストセル処理）は
付録~\ref{app:weno5} に収録する．
本文では概念的な役割と CCD/散逸フィルタとの関係のみを示す．

\begin{tcolorbox}[resultbox, title={本稿の空間離散化スキームの役割分担}]
\label{box:scheme_roles}
\begin{center}
\renewcommand{\arraystretch}{1.3}
\begin{tabular}{llll}
\hline
対象 & スキーム & 精度 & 根拠 \\
\hline
CLS 移流（$\psi$） & WENO5（付録~\ref{app:weno5}） & $O(h^5)$ & 非線形安定化・保存形 \\
曲率・速度微分 & CCD（§\ref{sec:CCD}） & $O(h^6)$ & 高精度・コンパクト \\
高周波ノイズ除去 & 散逸的 CCD（§\ref{sec:dissipative_ccd}） & — & CCD と同一 3 点ステンシル \\
圧力 Poisson & CCD-Poisson（§\ref{sec:pressure}） & $O(h^6)$ & Balanced-Force 整合性 \\
\hline
\end{tabular}
\end{center}
\end{tcolorbox}
